% =====================================================================
%  Appendix B -- Concordance.
%  The audit instrument for the "nothing omitted" claim: every heading of
%  study/tier0.md mapped to its location in this report, plus the tagged
%  equations and the three relocated self-check items.
% =====================================================================
\chapter{Concordance with the source document}
\label{app:concordance}

Every heading of \code{study/tier0.md} and where it landed in this report. The
source ordering is Part~0, I, II, \ldots, VII; this report's ordering is the
thematic one described in the front matter. Section numbers in the right column
are live cross-references, so they stay correct if the report is re-ordered
again.

\section{Headings}

\begingroup
\small
\setlength{\LTleft}{0pt}\setlength{\LTright}{0pt}
\begin{longtable}{@{}p{3.0cm} p{7.3cm} >{\raggedright\arraybackslash}p{5.2cm}@{}}
\toprule
\textbf{Source \S} & \textbf{Source heading} & \textbf{This report} \\
\midrule
\endfirsthead
\toprule
\textbf{Source \S} & \textbf{Source heading} & \textbf{This report} \\
\midrule
\endhead
\midrule
\multicolumn{3}{r@{}}{\footnotesize\emph{continued on next page}} \\
\endfoot
\bottomrule
\endlastfoot

\multicolumn{3}{@{}l}{\textbf{Front matter}} \\
--- & Title blurb & Title page \\
--- & Status of this document & \S\ref{sec:status} \\
--- & How to study each node & \S\ref{sec:howtostudy}, Fig.~\ref{fig:studyloop} \\
--- & Contents (table) & Reading map, \S\ref{sec:readingmap} \\
--- & Contents (full section list) & Appendix~\ref{app:sectionmap} \\
\addlinespace

\multicolumn{3}{@{}l}{\textbf{Part 0 --- Physical foundations}} \\
0 & Part preamble and prerequisite table & Part~\ref{part:substrate} opening \\
0.1   & The thermodynamic state of Si-burning matter & \S\ref{sec:thermostate} \\
0.1.1 & The four pressure components & \S\ref{sec:pressurecomponents} \\
0.1.2 & Worked example --- box centre & \S\ref{sec:boxcentre} \\
0.1.3 & Degeneracy across the box & \S\ref{sec:degeneracy} \\
0.1.4 & The Coulomb coupling parameter $\Gamma$ & \S\ref{sec:gamma} \\
0.1.5 & The electron-capture threshold & \S\ref{sec:ecthreshold} \\
0.1.6 & Self-check & after \S\ref{sec:ecthreshold} \\
0.2   & Statistical mechanics you actually need & \S\ref{sec:statmech} \\
0.2.1 & Chemical potential and the equilibrium condition &
        \S\ref{sec:chempotential} \\
0.2.2 & Nuclear partition functions & \S\ref{sec:partitionfunctions} \\
0.2.3 & Detailed balance from time reversal & \S\ref{sec:detailedbalance} \\
0.2.4 & Self-check & after \S\ref{sec:detailedbalance} \\
0.3   & Nuclear structure minimum & \S\ref{sec:nuclearstructure} \\
0.3.1 & The mass-excess convention & \S\ref{sec:massexcess} \\
0.3.2 & Magic numbers and \nuc{Ni}{56} & \S\ref{sec:magicnumbers} \\
0.3.3 & Where the rates actually come from & \S\ref{sec:ratesources} \\
0.3.4 & Self-check & after \S\ref{sec:ratesources} \\
0.4   & Neutrinos & \S\ref{sec:neutrinos} \\
0.4.1 & Two distinct neutrino channels & \S\ref{sec:nuchannels} \\
0.4.2 & Why neutrinos free-stream & \S\ref{sec:freestreaming} \\
0.4.3 & $\langle E_\nu\rangle$ --- an open item & \S\ref{sec:enu} \\
0.4.4 & Self-check & after \S\ref{sec:enu} \\
0.5   & Reaction nomenclature and REACLIB chapters & \S\ref{sec:nomenclature} \\
0.5   & Channel notation & \S\ref{sec:channelnotation} \\
0.5   & REACLIB chapters & \S\ref{sec:reaclibchapters},
        Fig.~\ref{fig:appendixb} \\
0.5   & In the code & repo-pointer box, \S\ref{sec:reaclibchapters} \\
0.6   & The topology of silicon burning & \S\ref{sec:topology} \\
0.6   & The $\alpha$ ladder & \S\ref{sec:alphaladder},
        Fig.~\ref{fig:alphaladder} \\
0.6   & Two clusters and the bridges & \S\ref{sec:clusters},
        Fig.~\ref{fig:clusters} \\
0.6   & A note on energy accounting & \S\ref{sec:energyaccounting} \\
0.7   & The timescale hierarchy & \S\ref{sec:timescales} \\
0.7   & The species destruction timescale & \S\ref{sec:tauspecies} \\
0.7   & The hierarchy, at box centre & \S\ref{sec:hierarchy} \\
0.7   & The separation that does not exist & \S\ref{sec:nogap} \\
0.7.1 & Self-check for \S0.5--0.7 &
        \S\ref{sec:selfcheck-network}, items 7--12 \\
\addlinespace

\multicolumn{3}{@{}l}{\textbf{Part I --- Why this project exists}} \\
I   & Part preamble & Part~\ref{part:context} opening \\
I.1 & The last day of a massive star & \S\ref{sec:lastday} \\
I.1 & The burning sequence & \S\ref{sec:burningsequence} \\
I.1 & Why the acceleration & \S\ref{sec:acceleration} \\
I.1 & Why this is the regime to emulate & \S\ref{sec:whyregime} \\
I.2 & What determines whether the star explodes & \S\ref{sec:explodability} \\
I.2 & Derivation --- $M_{\mathrm{ch}} \propto Y_e^2$ &
      \S\ref{sec:mchderivation} \\
I.2 & The thermal correction & \S\ref{sec:thermalcorrection} \\
I.2 & The collapse trigger & \S\ref{sec:collapsetrigger} \\
I.2 & The observable end of the chain & \S\ref{sec:observableend} \\
I.3 & Why networks are the computational bottleneck & \S\ref{sec:bottleneck} \\
I.3 & The three costs, quantified & \S\ref{sec:threecosts} \\
I.4 & Softwired networks & \S\ref{sec:softwired} \\
I.4 & The actual shape of the trade-off & \S\ref{sec:networkshape},
      Fig.~\ref{fig:nzplane} \\
I.4 & The measured consequence & \S\ref{sec:measuredconsequence} \\
I.5 & Why machine-learning emulation & \S\ref{sec:whyml} \\
I.5 & The cost model & \S\ref{sec:costmodel} \\
I.5 & Framing it honestly & \S\ref{sec:framinghonestly} \\
I.6 & What breaks when an emulator does not conserve &
      \S\ref{sec:conservationfailure} \\
I.6 & Mode 1 --- accumulation & \S\ref{sec:accumulation} \\
I.6 & Mode 2 --- physical inconsistency &
      \S\ref{sec:physicalinconsistency} \\
I.6 & Why soft penalties are the wrong fix & \S\ref{sec:softpenalties} \\
I.7 & Why a graph neural network & \S\ref{sec:whygnn} \\
I.8 & Where conservation is allowed to live &
      \S\ref{sec:whereconservation} \\
I.8 & Target A & \S\ref{sec:targetA} \\
I.8 & Target B and the projector derivation & \S\ref{sec:targetB} \\
I.8 & The critical theorem and its failure mode &
      \S\ref{sec:criticaltheorem} \\
I.9 & Self-check for Part I & items 1--3 $\to$
      \S\ref{sec:selfcheck-context}; items 4--6 $\to$
      \S\ref{sec:selfcheck-computational} \\
\addlinespace

\multicolumn{3}{@{}l}{\textbf{Part II (S0) --- Species and abundances}} \\
II.1 & Nuclide bookkeeping & \S\ref{sec:bookkeeping} \\
II.2 & Binding energy and the iron peak & \S\ref{sec:bindingenergy} \\
II.2 & The subtlety to internalize now & \S\ref{sec:constrainedoptimum} \\
II.3 & From number density to molar abundance &
       \S\ref{sec:molarabundance} \\
II.4 & Why networks integrate $Y$ and not $X$ & \S\ref{sec:whyY} \\
II.4 & Argument 1 --- per-baryon invariance & \S\ref{sec:perbaryon} \\
II.4 & Argument 2 --- the constraint stays linear &
       \S\ref{sec:linearconstraint} \\
II.4 & The network ODE & \S\ref{sec:networkode} \\
II.5 & The electron fraction \Ye{} & \S\ref{sec:ye} \\
II.5 & Derivation & \S\ref{sec:yederivation} \\
II.5 & The neutron excess & \S\ref{sec:neutronexcess} \\
II.5 & What changes \Ye{} & \S\ref{sec:whatchangesye} \\
II.5 & Why Si burning drives \Ye{} down & \S\ref{sec:yedown} \\
II.5 & The lepton ledger & \S\ref{sec:leptonledger} \\
II.6 & The network as a bipartite graph & \S\ref{sec:bipartite} \\
II.7 & Code walkthrough & \S\ref{sec:codewalkthrough-s0} \\
II.8 & Self-check for S0 & items 1--3, 5, 6, 8 $\to$
       \S\ref{sec:selfcheck-network} items 1--6; items 4 and 7 relocated
       (see below) \\
\addlinespace

\multicolumn{3}{@{}l}{\textbf{Part III (S1) --- The dataset}} \\
III.1  & What one record is & \S\ref{sec:onerecord} \\
III.2  & The regime box & \S\ref{sec:regimebox} \\
III.3  & The $\dd t$ grid & \S\ref{sec:dtgrid} \\
III.3  & Two deeper traps & \S\ref{sec:dttraps} \\
III.4  & Sobol sampling & \S\ref{sec:sobol} \\
III.4  & Discrepancy & \S\ref{sec:discrepancy} \\
III.4  & How Sobol is actually built & \S\ref{sec:sobolbuild} \\
III.4  & The gap you must fill honestly & \S\ref{sec:sobolgap} \\
III.4  & Scrambling & \S\ref{sec:scrambling} \\
III.4  & The consequence that shapes the repository &
         \S\ref{sec:nonregenerable} \\
III.5  & The identity problem --- \code{state\_id} & \S\ref{sec:stateid} \\
III.5  & Do this estimate yourself & \S\ref{sec:collisionestimate} \\
III.6  & Label noise, floors, censoring & \S\ref{sec:labelnoise} \\
III.6  & float32 labels & \S\ref{sec:float32labels} \\
III.6  & The $1\times10^{-15}$ floor & \S\ref{sec:floor} \\
III.6  & Energy normalization & \S\ref{sec:energynorm} \\
III.7  & Leakage-safe splits & \S\ref{sec:splits} \\
III.8  & The missing rows & \S\ref{sec:missingrows} \\
III.9  & Stratification & \S\ref{sec:stratification} \\
III.10 & Code reading order for S1 & \S\ref{sec:readingorder-s1},
         Fig.~\ref{fig:s1reading} \\
III.11 & Self-check for S1 & \S\ref{sec:selfcheck-dataset} \\
\addlinespace

\multicolumn{3}{@{}l}{\textbf{Part IV --- The modelling frame}} \\
IV   & Part preamble & Part~\ref{part:frame} opening \\
IV.1 & Operator splitting & \S\ref{sec:splitting} \\
IV.1 & The setup & \S\ref{sec:splitsetup} \\
IV.1 & Why this is the foundational assumption &
       \S\ref{sec:whyfoundational} \\
IV.1 & Consequence 1 --- constraint on $\Delta t$ &
       \S\ref{sec:splitconsequence1} \\
IV.1 & Consequence 2 --- inherited splitting error &
       \S\ref{sec:splitconsequence2} \\
IV.1 & Consequence 3 --- tractability &
       \S\ref{sec:splitconsequence3} \\
IV.2 & The one-zone approximation & \S\ref{sec:onezone} \\
IV.2 & What bbq is & \S\ref{sec:whatisbbq} \\
IV.2 & The subtlety nobody states & \S\ref{sec:convectivesubtlety} \\
IV.2 & The deepest structural insight &
       \S\ref{sec:reachablemanifold}, Fig.~\ref{fig:reachable} \\
IV.3 & The semigroup structure and rollout error &
       \S\ref{sec:rollouterror} \\
IV.3 & The semigroup property & \S\ref{sec:semigroup} \\
IV.3 & Rollout error propagation &
       \S\ref{sec:accumulationtrichotomy} \\
IV.4 & Self-check for Part IV & \S\ref{sec:selfcheck-frame} \\
\addlinespace

\multicolumn{3}{@{}l}{\textbf{Part V --- Numerical analysis}} \\
V.1 & Catastrophic cancellation & \S\ref{sec:cancellation} \\
V.1 & The derivation & \S\ref{sec:cancellationderivation} \\
V.1 & The numbers that follow & \S\ref{sec:kappanumbers} \\
V.1 & Why this reframes the kill-test & \S\ref{sec:killtestreframe} \\
V.2 & Conditioning and the rank-revealing subtlety &
      \S\ref{sec:conditioning} \\
V.3 & Stiffness, defined properly & \S\ref{sec:stiffness} \\
V.3 & Stiffness and QSE are the same phenomenon &
      \S\ref{sec:stiffnessqse} \\
V.4 & Floating point, concretely & \S\ref{sec:floatingpoint} \\
V.5 & Newton's method, damping, globalization & \S\ref{sec:newton} \\
V.5 & The method & \S\ref{sec:newtonmethod} \\
V.5 & Why it fails, specifically here & \S\ref{sec:newtonfails} \\
V.5 & Globalization & \S\ref{sec:globalization} \\
V.6 & ODE error control & \S\ref{sec:errorcontrol} \\
V.6 & Local versus global error & \S\ref{sec:localglobal} \\
V.6 & Why atol matters more than usual & \S\ref{sec:atol} \\
V.6 & Dense output and \code{t\_eval} & \S\ref{sec:denseoutput} \\
V.6 & BDF versus Radau & \S\ref{sec:bdfradau} \\
V.7 & Self-check for Part V & \S\ref{sec:selfcheck-numerics} \\
\addlinespace

\multicolumn{3}{@{}l}{\textbf{Part VI --- ML prerequisites and context}} \\
VI.1 & What kind of learning problem this is &
       \S\ref{sec:learningproblem} \\
VI.2 & Message passing, concretely & \S\ref{sec:messagepassing} \\
VI.3 & Feature and loss design & \S\ref{sec:featureloss} \\
VI.3 & Inputs first --- the legal transform &
       \S\ref{sec:inputtransform} \\
VI.3 & Then losses & \S\ref{sec:losses} \\
VI.4 & MESA and bbq, concretely & \S\ref{sec:mesabbq} \\
VI.5 & Prior art in network acceleration & \S\ref{sec:priorart} \\
VI.6 & The observational grounding &
       \textbf{moved} to \S\ref{sec:obsgrounding}
       (Part~\ref{part:context}) \\
VI.7 & Self-check for Part VI & items 1--5 $\to$
       \S\ref{sec:selfcheck-ml}; item 6 relocated (see below) \\
\addlinespace

\multicolumn{3}{@{}l}{\textbf{Part VII --- Working practice}} \\
VII   & Part preamble & Part~\ref{part:practice} opening \\
VII.1 & Every number carries a provenance label &
        \S\ref{sec:provenance} \\
VII.2 & Thresholds must have instruments & \S\ref{sec:instruments} \\
VII.3 & Decisions carry their reversal condition &
        \S\ref{sec:switchconditions} \\
VII.4 & Invariants live in code that refuses & \S\ref{sec:guards} \\
VII.5 & Non-regenerable artifacts are persisted &
        \S\ref{sec:persistence} \\
VII.6 & Notebooks explore; scripts produce & \S\ref{sec:notebooks} \\
VII.7 & Self-check for Part VII & \S\ref{sec:selfcheck-practice} \\
\addlinespace

\multicolumn{3}{@{}l}{\textbf{Closing matter}} \\
--- & Where Tier 0 hands off & \S\ref{sec:handoff} \\
--- & What Part 0 upgraded & ``What the substrate part upgraded'',
      \S\ref{sec:handoff} \\
--- & Threads that carry forward & \S\ref{sec:handoff} \\
--- & The three ideas to carry & \S\ref{sec:handoff} \\
--- & \textbf{Next:} Tier 1 trailer & end of \S\ref{sec:handoff} \\
--- & References (89 entries) & References \\

\end{longtable}
\endgroup

\section{Relocated self-check items}

Three of the source document's 73 self-check questions tested material that the
re-ordering moved. Each was relocated so that no question depends on a part the
reader has not yet reached. Nothing was dropped, and the count is unchanged.

\begin{center}
\small
\begin{tabularx}{\textwidth}{@{}l R{0.95} l R{1.05}@{}}
\toprule
\textbf{Source item} & \textbf{Question} & \textbf{Now in} &
\textbf{Reason} \\
\midrule
\S II.8 item 4 &
  \msn{151}-only $\beta$-decay chains; \msn{80}'s \Ye{} bias &
  \S\ref{sec:selfcheck-computational}, item \ref{sc:mesadiff} &
  requires the network diff of \S\ref{sec:softwired}, now later \\
\S II.8 item 7 &
  $\sum A u = 0$ vs $\sum A\,\dd Y = 0$ under $\asinh$ &
  \S\ref{sec:selfcheck-computational}, item 5 &
  requires the projector theorem of \S\ref{sec:criticaltheorem}, now later \\
\S VI.7 item 6 &
  1\% \Ye{} error $\to$ a changed observable &
  \S\ref{sec:selfcheck-context}, item 4 &
  tests \S VI.6, which now opens the report as \S\ref{sec:obsgrounding} \\
\bottomrule
\end{tabularx}
\end{center}

\section{Numbered equations}

The source tagged 21 equations. All are reproduced with semantic labels; two
load-bearing untagged equations were additionally numbered, marked \textbf{new}.

\begingroup
\small
\begin{longtable}{@{}l p{7.4cm} l@{}}
\toprule
\textbf{Source tag} & \textbf{Equation} & \textbf{This report} \\
\midrule
\endfirsthead
\toprule
\textbf{Source tag} & \textbf{Equation} & \textbf{This report} \\
\midrule
\endhead
\bottomrule
\endfoot
\bottomrule
\endlastfoot
(0.1)  & Chemical equilibrium, $\sum_i \nu_i\mu_i = 0$ &
         \eqref{eq:chem-equilibrium} \\
(0.2)  & NSE condition, $\mu(Z,N) = Z\mu_p + N\mu_n$ & \eqref{eq:nse} \\
(0.3)  & QSE cluster offset, $\mu_i = Z_i\mu_p + N_i\mu_n + u_G$ &
         \eqref{eq:qse-offset} \\
(0.4)  & Internal partition function $G(T)$ & \eqref{eq:partition-fn} \\
(0.5)  & Mass excess $\Delta(A,Z)$ & \eqref{eq:mass-excess} \\
(0.6)  & Binding energy from mass excesses & \eqref{eq:binding} \\
(0.7)  & $Q$-value from mass excesses & \eqref{eq:qvalue} \\
(0.8)  & Species destruction timescale $\tau_i$ & \eqref{eq:tau-species} \\
---    & $\kappa_r$ as the detailed-balance test \textbf{(new number)} &
         \eqref{eq:kappa-detailed-balance} \\
(1)    & Molar abundance $Y_i = X_i/A_i$ & \eqref{eq:molar-abundance} \\
(2)    & Baryon conservation $\sum A_i Y_i = 1$ &
         \eqref{eq:baryon-conservation} \\
(3)    & The network ODE $\dd\mathbf{Y}/\dd t = \nu\mathbf{R}$ &
         \eqref{eq:network-ode} \\
(4)    & Electron fraction \Ye{} & \eqref{eq:ye} \\
(5)    & Neutron excess $\eta = 1 - 2\Ye$ &
         \eqref{eq:neutron-excess} \\
---    & Ultra-relativistic degenerate EOS \textbf{(new number)} &
         \eqref{eq:urel-eos} \\
---    & $M_{\mathrm{ch}} \approx 5.83\,Y_e^2\,\Msun$ \textbf{(new number)} &
         \eqref{eq:mch} \\
---    & Projector $P = I - C\Tr(CC\Tr)^{-1}C$ \textbf{(new number)} &
         \eqref{eq:projector} \\
(IV.1) & Semigroup property & \eqref{eq:semigroup} \\
(IV.2) & Rollout error recursion & \eqref{eq:rollout-recursion} \\
(V.1)  & Cancellation error bound & \eqref{eq:cancellation-bound} \\
(V.2)  & $\kappa$ as reciprocal condition number &
         \eqref{eq:kappa-condition} \\
(V.3)  & Stiffness ratio $S$ & \eqref{eq:stiffness-ratio} \\
(V.4)  & Newton iteration & \eqref{eq:newton} \\
(V.5)  & Local error tolerance & \eqref{eq:error-tolerance} \\
(VI.1) & Message-passing update & \eqref{eq:message-passing} \\
\end{longtable}
\endgroup

\section{Diagrams}

The source document carried twelve fenced blocks. Seven were redrawn as TikZ
figures; five are literal listings (file paths, bin edges, a Fortran inlist, a
pair of quoted strings) and are set as such.

\begin{center}
\small
\begin{tabularx}{\textwidth}{@{}R{1.000} l@{}}
\toprule
\textbf{Source block} & \textbf{This report} \\
\midrule
Four-step study loop & Fig.~\ref{fig:studyloop} \\
Full section list & Appendix~\ref{app:sectionmap}, listing \\
Appendix-B config rows with annotations & Fig.~\ref{fig:appendixb} \\
The $\alpha$ ladder & Fig.~\ref{fig:alphaladder} \\
Two clusters and the bridges & Fig.~\ref{fig:clusters} \\
\msn{80} vs \msn{151} on the $(N,Z)$ plane & Fig.~\ref{fig:nzplane} \\
$\dd t$-grid file pointers & table, \S\ref{sec:dtgrid} \\
Stratification bin edges & table, \S\ref{sec:stratification} \\
S1 code reading order & Fig.~\ref{fig:s1reading} \\
Campaign inlist excerpt & listing, \S\ref{sec:whyfoundational} \\
Reachable manifold in composition space & Fig.~\ref{fig:reachable} \\
ADR incomplete/complete pair & table, \S\ref{sec:switchconditions} \\
\bottomrule
\end{tabularx}
\end{center}
